% 第 3 题频率分布直方图：按固定源 PDF 矢量路径实测重绘（1 pt → 0.03515 cm）。
% 原点取轴交点（源页 (132.48, 458.06) pt），y 向上。组距 2.5 小时 = 21.26 pt（横轴 1 组 = 21.26）。
% 纵轴刻度自下而上 0.02/0.04/0.08/0.10/0.16（无 0 标签，非单调步长）；柱高依次 0.02/0.10/0.16/0.08/0.04
% （实测 14.17/70.85/113.35/56.67/28.33 pt，与刻度完全自洽）。
% 5 条水平虚线自纵轴右行、止于对应柱左边缘（横穿其左侧各柱）；纵轴标签「频率／组距」为竖排分数式。
\begin{tikzpicture}[>=stealth,every node/.style={font=\small},x=0.03515cm,y=0.03515cm]
% 坐标轴（纵轴至 138.10、横轴至 166.43，箭头实心）
\draw[->] (0,0)--(0,141.7);
% 横轴：原点右侧带折断符号（源横轴基线上锯齿 x 136.73–149.48、振幅 ±2.13 pt，照源复刻；
% 初版遗漏，2026-10-09 由 2016 山东文（b20）转录对照回查发现并补齐）
\draw (0,0)--(4.25,0);
\draw (4.25,0)--(5.32,2.13)--(7.44,-2.13)--(9.57,2.13)--(11.69,-2.13)--(13.82,2.13)--(15.94,-2.13)--(17.01,0);
\draw[->] (17.01,0)--(170.0,0);
% 柱（高 0.02/0.10/0.16/0.08/0.04，组距 21.26）
\foreach \i/\h in {0/14.17, 1/70.85, 2/113.35, 3/56.67, 4/28.33} {
  \draw ({21.25+21.26*\i},0) rectangle ({21.25+21.26*(\i+1)},{\h});
}
% 水平虚线引导线（自纵轴至对应柱左缘）
\draw[dash pattern=on 2.99pt off 2.99pt] (0,14.17)--(21.25,14.17);
\draw[dash pattern=on 2.99pt off 2.99pt] (0,70.85)--(42.51,70.85);
\draw[dash pattern=on 2.99pt off 2.99pt] (0,113.35)--(63.77,113.35);
\draw[dash pattern=on 2.99pt off 2.99pt] (0,56.67)--(85.03,56.67);
\draw[dash pattern=on 2.99pt off 2.99pt] (0,28.33)--(106.29,28.33);
% 纵轴刻度标签（虚线高度处，纵轴左侧）
\node[left] at (-2.4,14.17) {\(0.02\)};
\node[left] at (-2.4,28.33) {\(0.04\)};
\node[left] at (-2.4,56.67) {\(0.08\)};
\node[left] at (-2.4,70.85) {\(0.10\)};
\node[left] at (-2.4,113.35) {\(0.16\)};
% 横轴刻度标签（组边界，居中于刻度）
\node[below] at (0,-2.6) {\(0\)};
\foreach \t/\x in {17.5/21.25, 20/42.51, 22.5/63.77, 25/85.03, 27.5/106.29, 30/127.55} {
  \node[below] at (\x,-2.6) {\(\t\)};
}
% 纵轴标签「频率／组距」竖排分数式（源页为分式排法：频率在上、横线、组距在下）
\node[anchor=west,inner sep=0] at (4.89,145.5) {频率};
\draw (4.89,140.99)--(25.89,140.99);
\node[anchor=west,inner sep=0] at (4.89,131.2) {组距};
% 横轴标签「自习时间/小时」（横轴上方、右端；基线在横轴上方 6.32）
\node[anchor=base west,inner sep=0] at (135.9,6.32) {自习时间/小时};
\end{tikzpicture}
